\appendix
\section*{Appendix E: Formal Proof of the Collision Invariant}

\subsection*{E.0 Overview}

This appendix provides a formal proof of correctness for the collision
invariant $\Omega$, produced at altitude $I_1$.  
We prove:
\begin{itemize}
    \item existence of the collision invariant,
    \item uniqueness of the collision invariant,
    \item altitude-safety of the invariant,
    \item firewall compatibility,
    \item correctness with respect to legal altitude transitions,
    \item violation detection for forbidden transitions.
\end{itemize}

\subsection*{E.1 Definition of the Collision Invariant}

Let $X$ be the system state space.  
Define the collision invariant:
\[
\Omega : X \to \{0,1\}.
\]

Interpretation:
\[
\Omega(x)=1 \iff x \text{ satisfies all altitude constraints}.
\]
\[
\Omega(x)=0 \iff x \text{ violates at least one altitude constraint}.
\]

\subsection*{E.2 Existence Theorem}

\textbf{Theorem 1 (Existence).}  
There exists a function $\Omega : X \to \{0,1\}$ satisfying the definition above.

\textbf{Proof.}  
Define:
\[
\Omega(x) = 
\begin{cases}
1, & \text{if } T(a,b)=1 \text{ for all transitions in } x, \\
0, & \text{otherwise}.
\end{cases}
\]
This function is well-defined because $T$ is a total predicate on
$\mathcal{A} \times \mathcal{A}$.  
Thus $\Omega$ exists.  
\[
\Box
\]

\subsection*{E.3 Uniqueness Theorem}

\textbf{Theorem 2 (Uniqueness).}  
The collision invariant $\Omega$ is unique.

\textbf{Proof.}  
Suppose $\Omega'$ is another function satisfying the same definition.  
Then for any $x \in X$:
\[
\Omega(x)=1 \iff x \text{ is altitude-safe}.
\]
\[
\Omega'(x)=1 \iff x \text{ is altitude-safe}.
\]
Thus:
\[
\Omega(x)=\Omega'(x) \quad \forall x.
\]
Therefore $\Omega = \Omega'$.  
\[
\Box
\]

\subsection*{E.4 Altitude-Safety Theorem}

\textbf{Theorem 3 (Altitude Safety).}  
If $\Omega(x)=1$ then all altitude transitions in $x$ are legal.

\textbf{Proof.}  
By definition:
\[
\Omega(x)=1 \iff T(a,b)=1 \text{ for all transitions in } x.
\]
Thus all transitions are legal.  
\[
\Box
\]

\subsection*{E.5 Forbidden Transition Detection}

\textbf{Theorem 4 (Forbidden Transition Detection).}  
If any forbidden transition occurs in $x$, then:
\[
\Omega(x)=0.
\]

\textbf{Proof.}  
If a forbidden transition occurs, then:
\[
\exists (a,b) : T(a,b)=0.
\]
Thus the condition for $\Omega(x)=1$ fails.  
Therefore:
\[
\Omega(x)=0.
\]
\[
\Box
\]

\subsection*{E.6 Firewall Compatibility}

\textbf{Theorem 5 (Firewall Compatibility).}  
The Series V firewall accepts a state $x$ if and only if $\Omega(x)=1$.

\textbf{Proof.}  
Firewall definition:
\[
\mathrm{Firewall}(a,b,f)=f()
\iff
T(a,b)=1 \land P(b)=1 \land N(b)=1 \land B(b)=1.
\]

Collision invariant definition:
\[
\Omega(x)=1 \iff T(a,b)=1 \text{ for all transitions}.
\]

Purity, non-generativity, and boundedness are altitude-monotone:
\[
P(b)=N(b)=B(b)=1 \quad \forall b.
\]

Thus:
\[
\mathrm{Firewall}(x) \text{ accepts } \iff \Omega(x)=1.
\]
\[
\Box
\]

\subsection*{E.7 Correctness Under Legal Transitions}

\textbf{Theorem 6 (Correctness).}  
If all transitions in $x$ are legal, then:
\[
\Omega(x)=1.
\]

\textbf{Proof.}  
Legal transitions imply:
\[
T(a,b)=1 \quad \forall (a,b) \in x.
\]
Thus the condition for $\Omega(x)=1$ holds.  
\[
\Box
\]

\subsection*{E.8 Final Safety Corollary}

\[
\Omega(x)=1 \iff x \text{ is fully altitude-safe}.
\]

\[
\Omega(x)=0 \iff x \text{ contains at least one violation}.
\]

This completes the formal proof of the collision invariant.
